\chapter*{绪论}
\addcontentsline{toc}{chapter}{绪论}
\markboth{绪论}{绪论}
\setcounter{section}{0}
\renewcommand{\thesection}{\arabic{section}}

数值线性代数又称矩阵计算（matrix computation）。
中心任务不是再证一遍存在唯一性，而是：针对矩阵问题的结构，设计\textbf{快速且可靠}的算法。

\section{数值线性代数的基本问题}

三类核心问题：
\begin{align}
Ax&=b,\\
\min_y\,\lVert Ay-b\rVert_2,\\
Ax&=\lambda x\quad(x\neq 0).
\end{align}
第一式输入非奇异 $A\in\mathbf{R}^{n\times n}$ 与 $b\in\mathbf{R}^{n}$，输出未知量 $x$。
第二式输入 $A\in\mathbf{R}^{m\times n}$、$b\in\mathbf{R}^{m}$，输出使残差欧氏长度最小的 $x$；该长度非负，一般不能把残差降到零。
第三式输出特征值 $\lambda$ 与特征向量 $x$（可差非零倍数）。
奇异值分解因应用广泛，现常列为第四类问题。

\section{研究数值方法的必要性}

理论上的闭式解往往不能当算法用。
Cramer 法则把 $x_i$ 写成两个行列式之比，再配 Laplace 展开，解 $n$ 阶组的乘法次数大于 $(n+1)!$。
在每秒 $10^{10}$ 次运算的机器上取 $n=25$，约需 $13$ 亿年；同一问题用消元法不到一秒。
$n!$ 是运算次数，不是解的某个分量——它说明「公式漂亮」与「可计算」不是一回事。
特征值的 Jordan 分解能读出重数，但对扰动不稳定；实际计算改用 Schur 分解。

\section{矩阵分解是设计算法的主要技巧}

全书的设计思想可以压成一句：
把一般矩阵问题化成一个或几个\textbf{容易求解的特殊问题}，而完成这一转化的主要技巧是矩阵分解。

\bookfig{fig-intro-triangular.png}{绪论第4页}{下三角组与上三角组可以前代、回代直接解出。}{$L$：下三角；$U$：上三角；$x,y$：未知量；$b,c$：右端}

一般的 $Ax=b$ 先分解
\[
PA=LU,
\]
再解 $Ly=Pb$ 与 $Ux=y$。
$P$ 是排列方阵，$L$ 下三角，$U$ 上三角；$y$ 只是中间量，$x$ 才是原未知量。
于是「解线性组」变成「如何分解」加上两次 $O(n^2)$ 的三角形求解。这是第一章的主线。

\section{敏度分析与误差分析}

计算结果很少等于真值。误差来自数据本身，也来自计算过程。
要问「差多少」，必须把\textbf{问题}和\textbf{算法}分开看。

敏度分析：数据微扰后，真解变多少。在相对扰动很小的前提下，
\begin{equation}
\frac{\lvert f(x+\delta x)-f(x)\rvert}{\lvert f(x)\rvert}
\le c(x)\frac{\lvert\delta x\rvert}{\lvert x\rvert}.
\end{equation}
$c(x)$ 是 $f$ 在 $x$ 处的条件数（condition number），无量纲。
$c(x)$ 大称病态，小称良态——这是问题的固有属性，与算法无关。
可微时 $c(x)\approx\lvert f'(x)\rvert\,\lvert x\rvert/\lvert f(x)\rvert$。

误差分析（后向）：把算出的 $\hat y$ 看成精确作用在受扰数据上的结果，
\[
\hat y=f(x+\delta x),\qquad \lvert\delta x\rvert\le\lvert x\rvert\,\varepsilon.
\]
$\varepsilon$ 小称算法数值稳定——这是算法的固有属性，与问题是否病态无关。

合在一起：
\begin{equation}
\frac{\lvert\hat y-f(x)\rvert}{\lvert f(x)\rvert}\le c(x)\,\varepsilon.
\end{equation}
相对误差上界是条件数与后向扰动水平之积。
只有对良态问题使用稳定算法，才可期望可靠结果。

\section{算法复杂性与收敛速度}

直接法：无舍入时有限步得到精确解，快慢主要看运算量（现计全部四则，并略去低阶，如 $3n^3+20n^2$ 称作 $3n^3$）。
运算量不能单独决定快慢：数据传输往往比算术更贵。

迭代法：有限步得不到精确解，还要看
\[
\lVert x_k-x\rVert\le c\lVert x_{k-1}-x\rVert
\quad(0<c<1)
\]
（线性收敛）还是
\[
\lVert x_k-x\rVert\le c\lVert x_{k-1}-x\rVert^2
\]
（平方/二次收敛）。$\lVert x_k-x\rVert$ 是第 $k$ 步误差的长度；$c$ 越小越好。

\section{算法的软件实现与现行数值线性代数软件包}

可靠实现需要算法理论、机器结构与计算经验，不是把公式翻译成循环。
基本算法主要收在 LAPACK 与 MATLAB 中；后者便于试用，一般慢于直接调 LAPACK。

\section{符号说明}

矩阵用大写 $A,B,\Lambda$，元素 $a_{ij}$，向量 $x,y$，单位阵 $I$，其第 $i$ 列为 $e_i$。
排列方阵 $P=[e_{i_1},\ldots,e_{i_n}]$。
$\operatorname{tr}A$、$\operatorname{rank}A$、$\det A$、$A^{-1}$、$A^{\mathsf T}$、$A^{-\mathsf T}=(A^{-1})^{\mathsf T}$。
算法下标沿用 $i:j$、$A(i,:)$、$A(k:l,p:q)$，分别表示整数段、一行、一个子块。

\renewcommand{\thesection}{\thechapter.\arabic{section}}
